\section{Оценивание точности метода извлечения антропометрических признаков}
В этой части представлен результат извлечения антропометрических признаков на основе алгоритма компьютерного зрения из $50$ изображений. Извлечение ключевых точек (меток - опорных точек) состоит из $2$ главных этапов:

\begin{itemize}
	\item \textbf{Шаг 1:} Предварительная обработка изображения, преобразование исходного цветного изображения в изображение в оттенках серого, фильтрация шумов, разряжение пикселей, эквализация гистограммы для получения лучшего качества изображения;
	\item \textbf{Шаг 2: } Извлечение точек признаков с использованием алгоритма вычитания фона, обнаружения человеческого лица, сегментации изображения, поиска опорных точек для вычисления антророметрических признаков.
\end{itemize}

Для оценки эффективности процесса извлечения признаков применяются средняя абсолютная ошибка - $MAPE$:
\begin{equation}\label{eq26}
M=\frac{1}{n}\sum^n_{i=1}\left|\frac{A_t-F_t}{A_t}\right|,
\end{equation}
Где:

\begin{itemize}
	\item $A_t$: Результат измерений, рассчитанных вручную;
	\item $F_t$: Результат извлечения антропометрических признаков.
\end{itemize}
Эксперимент был использован для проверки погрешности размеров человеческих частей на основе извлечения количества признаков. На основании результатов для калибровки точности алгоритма.

\textbf{Случай 1}: Результат извлечения 24 опорных точек на контуре человеческого тела. Иллюстрация расположения опорных точек (рис. \ref{img15}): 
\begin{figure}[ht!]
\centering
\includegraphics [scale=0.8] {images/h15.png}
\begin{center}
%\captionsetup{justification=justified, labelsep=period}
\caption{Результат расположения 24 опорных точек \cite{long1,long2}.} \label{img15}
\end{center}
\end{figure}

\begin{table}[b!]%
\begin{center}
\caption{Результат извлечения антропометрических признаков и погрешности системы компьютерного зрения \cite{long1,long2}.}\label{tab1}
  \begin{tabular}{|c|c|c|c|c|c|c|}
    \hline
    \multirow{2}{*}{Антропометрические} & {Ручной} & \multicolumn{3}{c}{Результаты системы} & {$\sum^n_{i=1}\left|\frac{A_t-F_t}{A_t}\right|$} &{$M$} \\
      признаки& метод &1 &2 &3& & \\
    \hline
Груди &95	&93.65	&94.11	&93.32	&0.04126	&0.01375\\
\hline
Талия               &79	&76.06	&78.80	&80.15	&0.05430	&0.01810\\
\hline
Бедра               &96.5	&93.07	&95.19	&95.78	&0.05658	&0.01886\\
\hline
Длина рук           &53	  &52.45	&54.32	&54.72	&0.06774	&0.02258\\
\hline
Обхват бицепса      &28	 &25.68	  &27.09	&27.60	&0.12964	&0.04321\\
\hline
Обхват шеи          &37	 &35.71	  &36.67	&36.79	&0.04946	&0.01649\\
\hline
Длина спины         &40	 &38.50	  &40.11	&40.01	&0.04050	&0.01350\\
\hline
Длина плеча       &36	 &35.86	  &36.19	&35.50	&0.02306	&0.00769\\
\hline
Ширина плеча         &14	 &13.91	  &13.56	&14.74	&0.09071	&0.03024\\
\hline
Высота груди        &18	 &17.78	  &17.60	&17.86	&0.04222	&0.01407\\
\hline
  \end{tabular}
\end{center}
\end{table}%\vspace{10mm}

\textbf{Случай 2:} Применение метода калибровки для улучшения точности извлечения признаков (рис. \ref{img152}).
\begin{figure}[ht!]
\centering
\includegraphics [scale=0.8] {images/h152.png}
\begin{center}
%\captionsetup{justification=justified, labelsep=period}
\caption{Результат расположения 28 опорных точек  \cite{long1,long2}.} \label{img152}
\end{center}
\end{figure}
\begin{table}[b!]%
\begin{center}
\caption{Результат извлечения антропометрических признаков и погрешности системы компьютерного зрения после применения калибровки извлеченных точек \cite{long1,long2}.}.\label{tab3}
  \begin{tabular}{|c|c|c|c|c|c|c|}
    \hline
    \multirow{2}{*}{Антропометрические} & {Ручной} & \multicolumn{3}{c}{Результаты системы} & {$\sum^n_{i=1}\left|\frac{A_t-F_t}{A_t}\right|$} &{$M$} \\
      признаки& метод  &1 &2 &3 & & \\
    \hline
Груди &95	&95.03	&95.00	&95.02	&0.00053	&0.00018 \\
\hline 
Талия             &79	&78.46	&79.05	&78.89	&0.00557	&0.00186\\

\hline
Бедра               &96.5	&97.29	&96.11	&96.4	&0.00010	&0.00003\\

\hline
Длина рук           &53	&52.93		&52.1	&52.93	&0.03531	&0.01177\\

\hline
Обхват бицепса      &28	&28.24	&28.03	&28.07	&0.01071	&0.00357\\

\hline
Обхват шеи         &37	&38.35	&38.13	&38.05	&0.09672	&0.03224\\

\hline
Длина спины         &40	&38.35	&39.96	&39.81	&0.04325	&0.01442\\

\hline
Длина плеча       &36	&36.01	&36.17	&36.1	&0.00971	&0.00324\\

\hline
Ширина плеча        &14	&13.97	&14.02	&14.02	&0.00071	&0.00024\\

\hline
Высота груди       &18	&18.05	&18.11	&18.01	&0.01500	&0.00500\\

\hline
  \end{tabular}
\end{center}
\end{table}%\vspace{10mm}

\begin{figure}[ht!]
\centering
\includegraphics [scale=0.7] {images/h16.png}
\begin{center}
%\captionsetup{justification=justified, labelsep=period}
\caption{Погрешность в обоих случаях в извлечении антропометрических признаков.} \label{img16}
\end{center}
\end{figure}

\begin{figure}[ht!]
\centering
\includegraphics [scale=0.7] {images/h18.png}
\begin{center}
%\captionsetup{justification=justified, labelsep=period}
\caption{Процент правильной классификации (\%)  в обоих случаях в извлечении антропометрических признаков.} \label{img18}
\end{center}
\end{figure}

Для вычичления технической погрешности измерений при извлечении антропометрических признаков используется метод $TEM$(technical error of measurement)\cite{Stanley1999}. В данной работе выполняется извлечение антропометрических признаков 100 раз для одного человека.

С помощью метода $TEM$:

	\begin{equation}\label{eq50}
TEM=\sqrt{\left(\left(\sum_1^N\left(\left(\sum_1^K M^2\right)-\left(\left(\sum_1^K M\right)^2\right)\right)\right)/N\left(K-1\right)\right)}
\end{equation}
Где

\begin{itemize}
	\item $N$ - количество объектов;
	\item $K$ - количество выполнения;
	\item $M$ - измерение.
\end{itemize}

Получена гистограмма технической погрешности измерений (рис. \ref{img45}).

\begin{figure}[ht!]
\centering
\includegraphics [scale=1] {images/h45.png}
\begin{center}
%\captionsetup{justification=justified, labelsep=period}
\caption{Результат анализа погрешности.} \label{img45}
\end{center}
\end{figure}

В (рис. \ref{img45}) показано, что среднее арифметичесное $mean = 29050.66752$, медиана $median = 29039.11819$ и скошенность $skewness= 0.423$. В этом распределении среднее арифметичесное и медиана приблизительны, скошенность $\in \left[-1;1\right]$. Поэтому это нормальное распределение (распределение Гаусса) (рис. \ref{img44}).

\begin{figure}[ht!]
\centering
\includegraphics [scale=1] {images/h44.png}
\begin{center}
%\captionsetup{justification=justified, labelsep=period}
\caption{Гистограмма технической погрешности измерений.} \label{img44}
\end{center}
\end{figure}



  